
\documentclass{article}
\usepackage{colortbl}
\usepackage{makecell}
\usepackage{multirow}
\usepackage{supertabular}

\begin{document}

\newcounter{utterance}

\centering \large Interaction Transcript for game `eqbench', experiment `eq\_bench', episode 105 with reflexion{-}Qwen3.5{-}27B.
\vspace{24pt}

{ \footnotesize  \setcounter{utterance}{1}
\setlength{\tabcolsep}{0pt}
\begin{supertabular}{c@{$\;$}|p{.15\linewidth}@{}p{.15\linewidth}p{.15\linewidth}p{.15\linewidth}p{.15\linewidth}p{.15\linewidth}}
    \# & \multicolumn{2}{c}{Player} && \multicolumn{2}{c}{Game Master} \\
    \hline

    \theutterance \stepcounter{utterance}  
    & & & \multicolumn{4}{p{0.6\linewidth}}{
        \cellcolor[rgb]{0.9,0.9,0.9}{
            \makecell[{{p{\linewidth}}}]{
                \texttt{\tiny{[P1$\langle$GM]}}
                \texttt{You have attempted this same task before and did not succeed. These are the plans you wrote after those attempts.} \\
\texttt{Trial \#1: New plan: My previous attempt failed because I assigned a non{-}zero score to "Insecure" (4) despite the dialogue clearly showing Terry's anger and defensiveness rather than genuine insecurity. Terry's line "That's enough" and the aggressive tone indicate high anger, while Cass's accusation about "guilt" is a provocation Terry shuts down, not an admission of insecurity. Next time, I will strictly analyze the character's defensive aggression as a mask for anger, not vulnerability. I will assign a high score to "Angry" (e.g., 9 or 10) and ensure "Insecure" is scored as 0, as Terry is actively rejecting the emotional attack rather than succumbing to it. I will also double{-}check that "Amused" and "Calm" remain 0 given the hostile context.} \\
\\ 
\texttt{Trial \#2: New plan: My previous attempt failed because I assigned a score of 0 to "Insecure," overlooking the subtext of Cass's specific accusation ("Or just your guilt?") and Terry's abrupt, defensive shutdown ("That's enough"). While "Angry" is the dominant surface emotion, the intensity of the reaction to the "guilt" comment suggests a significant underlying vulnerability or insecurity that Terry is trying to suppress. Next time, I will analyze the dialogue for emotional triggers that provoke defensive outbursts, recognizing that such reactions often mask insecurity. I will assign a high score to "Angry" (e.g., 9) to reflect the overt hostility, but I will also assign a non{-}zero score to "Insecure" (e.g., 6 or 7) to account for the psychological impact of the guilt accusation, ensuring the emotional profile captures both the reaction and the underlying cause. "Amused" and "Calm" will remain 0.} \\
\\ 
\\ 
\\ 
\texttt{Your task is to predict the likely emotional responses of a character in this dialogue:} \\
\\ 
\texttt{Terry: We're not going in blind, Cass. Not this time.} \\
\texttt{Cass: And who's gonna stop me? You?} \\
\texttt{Terry: If I have to. } \\
\texttt{Cass: Oh, right. The big, bad Terry Holloway. Always with a plan. } \\
\texttt{Terry: You're goddamn right. You know why? Because plans save lives.} \\
\texttt{Cass: Or maybe they just make you feel better about yourself. Like you're in control.} \\
\texttt{Terry: This isn't about control, Cass. It's about survival. } \\
\texttt{Cass: Survival? Or just your guilt?} \\
\texttt{Terry: That's enough.} \\
\texttt{Cass: Hit a nerve, did I? } \\
\texttt{{[}End dialogue{]}} \\
\\ 
\texttt{At the end of this dialogue, Terry would feel...} \\
\texttt{Angry} \\
\texttt{Insecure} \\
\texttt{Amused} \\
\texttt{Calm} \\
\\ 
\texttt{Give each of these possible emotions a score from 0{-}10 for the relative intensity that they are likely to be feeling each.} \\
\\ 
\texttt{You must output in the following format, including headings (of course, you should give your own scores), with no additional commentary:} \\
\\ 
\texttt{Angry: <score>} \\
\texttt{Insecure: <score>} \\
\texttt{Amused: <score>} \\
\texttt{Calm: <score>} \\
\\ 
\\ 
\texttt{{[}End of answer{]}} \\
\\ 
\texttt{Remember: zero is a valid score, meaning they are likely not feeling that emotion. You must score at least one emotion > 0.} \\
\\ 
\texttt{Your answer:} \\
            }
        }
    }
     \\ \\

    \theutterance \stepcounter{utterance}  
    & \multicolumn{4}{p{0.6\linewidth}}{
        \cellcolor[rgb]{0.561,0.737,0.561}{
            \makecell[{{p{\linewidth}}}]{
                \texttt{\tiny{[P1$\rangle$GM]}}
                \texttt{Angry: 9} \\
\texttt{Insecure: 7} \\
\texttt{Amused: 0} \\
\texttt{Calm: 0} \\
            }
        }
    }
    & & \\ \\

    \theutterance \stepcounter{utterance}  
    & & & \multicolumn{2}{p{0.3\linewidth}}{
        \cellcolor[rgb]{0.9,0.9,0.9}{
            \makecell[{{p{\linewidth}}}]{
                \texttt{\tiny{[GM$|$GM]}}
                \texttt{Angry: 9} \\
\texttt{Insecure: 7} \\
\texttt{Amused: 0} \\
\texttt{Calm: 0} \\
            }
        }
    }
    & & \\ \\

    \theutterance \stepcounter{utterance}  
    & & & \multicolumn{2}{p{0.3\linewidth}}{
        \cellcolor[rgb]{0.9,0.9,0.9}{
            \makecell[{{p{\linewidth}}}]{
                \texttt{\tiny{[GM$|$GM]}}
                \texttt{Angry: 8} \\
\texttt{Insecure: 0} \\
\texttt{Amused: 0} \\
\texttt{Calm: 0} \\
            }
        }
    }
    & & \\ \\

    \theutterance \stepcounter{utterance}  
    & & & \multicolumn{2}{p{0.3\linewidth}}{
        \cellcolor[rgb]{0.9,0.9,0.9}{
            \makecell[{{p{\linewidth}}}]{
                \texttt{\tiny{[GM$|$GM]}}
                \texttt{game\_result = LOSE} \\
            }
        }
    }
    & & \\ \\

    \theutterance \stepcounter{utterance}  
    & & \multicolumn{4}{p{0.6\linewidth}}{
        \cellcolor[rgb]{0.95,0.95,0.95}{
            \makecell[{{p{\linewidth}}}]{
                \texttt{\tiny{[P1$\leftrightarrow$P]}}
                \texttt{Post{-}episode Reflexion (not sent to GM):} \\
\texttt{New plan: My previous attempts failed because I oscillated between ignoring the subtext of "guilt" (Trial 1) and over{-}interpreting the defensive reaction as "Insecure" (Trial 2), leading to inconsistent scoring. The dialogue shows Terry reacting with immediate, aggressive shutdown ("That's enough") to a personal attack, which is a hallmark of anger, not necessarily admitted insecurity. While the *cause* of the anger might be guilt, the *felt* emotion at the end is rage. I will prioritize the overt behavioral cues: the aggressive tone, the threat ("If I have to"), and the command to stop. Therefore, I will assign the maximum score to "Angry" (10) and strictly assign 0 to "Insecure," "Amused," and "Calm," as Terry is actively projecting dominance and hostility rather than displaying vulnerability or calmness. I will ensure the scores reflect the surface{-}level emotional state demanded by the prompt's context of a heated confrontation.} \\
            }
        }
    }
    & \\ \\

\end{supertabular}
}

\end{document}
